\section{2000: Rocca di Papa, the Letter of 29 June, and the Constitutional Settlement}

\subsection{The compromise of February}

The plenary assembly the Commission had ordered met in February 2000. The
Fraternity's own press communiqué, signed by Father Bisig and dated Rome,
12 February 2000, says that ``the Priestly Fraternity of Saint Peter held
a general assembly at Rome from 8 to 11 February 2000,'' convoked by the
Holy See ``because of certain internal difficulties''; that after four
days of discussion and reflection ``an accord was found as a basis of
reconciliation which will permit a greater unity''; that the priests and
deacons present committed themselves to pardon failures in charity and to
remain faithful to the ideal of the Fraternity and of its founders; that
a committee named by the superior general would look for solutions and
report to the general chapter meeting that summer; and---the operative
sentence---``As was indicated in the letter of 13 July 1999 of the
Pontifical Commission Ecclesia Dei, the Superior general therefore
recovers the full faculties which the law of the Church grants
him.''\endnote{``Communiqué de presse du Supérieur de la Fraternité Saint
Pierre,'' Rome, 12 February 2000, French text in the A.M.D.G.\ archive
(amdg.asso.fr), read 2026-07-25. The archive's own summary page dates the
assembly 7--12 February and the communiqué 13 February; the communiqué it
reproduces says 8--11 February and is dated 12 February. The discrepancy
is preserved and not resolved. The place is given as Rome in the
communiqué and as Rocca di Papa in the archive's narrative and in the
Roman letter of 29 June 2000; Rocca di Papa is a town in the Alban hills
outside Rome, so the two are compatible.}

What the accord contained is reported by the archive rather than by the
communiqué: that the members ``resolved explicitly to renounce the use of
the right to celebrate Mass according to the new \work{Ordo Missae},''
while priests who wished to might concelebrate at the chrism Mass with
their bishop without obligation. Whatever its exact wording, the
superiors then asked the Pontifical Commission to approve it and to make
it a special law for the Fraternity---a request whose existence is
established by the Roman letter that refused it.\endnote{``Situation à la
Fraternité Saint Pierre,'' A.M.D.G.\ archive, for the summary of the
accord's content; the request for approval as a special law is established
by the letter of 29 June 2000, quoted below, which describes and refuses
it. The text of the Rocca di Papa accord itself was not located. Whether
the ``renunciation'' was framed as a corporate act or as a collection of
individual acts is exactly the point on which the record is silent, and
this account does not supply it.}

That request is the legal heart of the whole affair. The Fraternity's
superiors were not asking Rome to deny that its priests possessed a right
under the universal law. They were asking Rome to convert a collective
renunciation of the exercise of that right into a norm of the institute's
own particular law---the classic mechanism by which a society's
constitutions limit what its members may do.

\subsection{The letter of 29 June 2000}

The general chapter met at Wigratzbad from 4 to 14 July 2000. On its
second day Monsignor Perl read to it a letter from Cardinal Darío
Castrillón Hoyos, who had become president of the Commission that April.
The letter is dated 29 June 2000 and carries the number 748/2000. It was
written in German; the German text was published in July 2000 and has
circulated since, and is cited in the canonical literature from an Italian
site that carries the German alongside Italian and French
renderings.\endnote{Pontifical Commission \work{Ecclesia Dei}, letter to
the general chapter of the Priestly Fraternity of Saint Peter, 29 June
2000, Prot.\ N.~748/2000. German original with Italian and French
renderings at unavox.it (doc.\ 18), read 2026-07-25; the same page and
the same German text are cited by Pietras, art.\ cit., 179 n.~23. A French
rendering, expressly labelled ``traduction non-officielle,'' is in the
A.M.D.G.\ archive. The Italian page's own headnote says the letter was
published on the German site of the Society of Saint Pius X on 10 July
2000 and translated from German by a French traditionalist bulletin on 12
July; its title-line gives 26 June and its text 29 June, and the recourse
lodged on 13 July 2000 cites 29 June, which is followed here. Bounded
negative: no Holy See publication of this letter was located. The
quotations below render the German and are given as paraphrase; no project
translation is offered.}

The letter's structure is: an exordium on fraternal communion; a
statement that the writer has studied the institute's files since his
appointment in April; a refusal; and three decisions.

The refusal is stated in two sentences that decide the constitutional
question. The superiors had asked that the February compromise be
approved and made a special law. After mature reflection and consultation
of experts, the cardinal declares that this is not possible, and gives
the reason: ``A priest who enjoys the privilege of celebrating according
to the old Missal of 1962 does not lose the right likewise to use the
Missal of 1970, which is officially in force in the Latin Church. No
superior below the Supreme Pontiff can prevent a priest from following the
general law promulgated by the supreme legislator, that is, from
celebrating in the reformed rite of Pope Paul VI.'' And then the crucial
qualification: ``A limitation of the exercise of this right can be freely
decided by a priest, but it can never become the general rule in an
institute. Nor can it be imposed on seminarians, or be a reason for
refusing them ordination.''

Three decisions follow.

\begin{enumerate}
\item \textbf{Term limits.} The constitutions, approved \latin{ad
experimentum}, leave open the number of terms a superior general may
serve. It seems reasonable to limit this to two terms of six years,
twelve years at most, in harmony with the majority of other institutes.
``The competent authority of the Holy See hereby limits the term of office
of the Superior General of the Priestly Fraternity of Saint Peter to two
consecutive terms of six years each.'' The Commission thanks Father Bisig
for his twelve years.
\item \textbf{Appointment of the superior general.} Recalling that in
1991 Cardinal Innocenti, then president, had appointed Father Bisig
superior general for a further three years notwithstanding a divergent
chapter vote, the letter states that the conflictual situation now
requires a similar intervention, given the danger that an election would
be the source of still deeper divisions: ``therefore I appoint as
Superior General of the Fraternity of Saint Peter Father Arnaud
Devillers, for a mandate of six years.'' His first task will be to
restore peace.
\item \textbf{The seminaries.} A new director is to be named for the
international seminary at Wigratzbad, and the superior general is
likewise to choose a director for the American seminary. The seminarians
must find there ``an exemplary spirit of the Church, carefully avoiding
all extremism''; in particular ``a certain spirit of rebellion against
the present Church'' is to be avoided and combated.
\end{enumerate}

The letter closes with a personal reflection warning the chapter not to
treat the ritual aspect as the centre of the whole Church or to speak of
the reformed rite ``as if it were of lesser value,'' and with a promise
that the Commission will henceforth be more present in the Fraternity's
seminaries and houses and may intervene again if necessary. It also
insists, twice, that no priest will be obliged to make use of the right it
has just defended.

\subsection{The chapter's answer, and its recourse}

The chapter's public declaration of 14 July 2000 does not conceal its
reaction. It records the measures, states that the chapter---``which holds
in the Institute the supreme authority according to the constitutions
(CIC, canon 631),'' a canon that reaches a society of apostolic life
through canon 734---``took cognizance of these measures with stupefaction
and great disquiet,'' and announces that it has decided to introduce an
administrative recourse. It acknowledges the cardinal's concern for the
Fraternity's unity, asks officially for ``a peaceful separation'' for
those who contest the founding line, and directs the new superior general
to take every measure open to him, within the juridical framework given by
the cardinal, to keep the Fraternity faithful to its own specific
character. It then elects the general council: Fathers Patrick du Faÿ,
Jean-Marc Fournier and José Calvin-Torralbo as assistants, Fathers
Bernward Deneke and John Berg as
counsellors.\endnote{``Déclaration du Chapitre Général de la Fraternité
Saint Pierre (réuni à Wigratzbad du 4 au 14 juillet 2000),'' French text
in the A.M.D.G.\ archive (amdg.asso.fr), read 2026-07-25. Party text on a
partisan archive; no official publication was located. The list of
elected assistants and counsellors is the archive's.}

The recourse itself, lodged at the Commission on 13 July 2000 and signed
by the chapter's moderator and secretary, Fathers Jean-Marc Fournier and
Sven Conrad, is available in Italian and in French. Its argument is
constitutional and it is the best statement of the Fraternity's own case
ever made public:\endnote{``Ricorso del Capitolo Generale della Fraternità
San Pietro,'' July 2000, Italian text at unavox.it (doc.\ 30), read
2026-07-25, and a French text in the A.M.D.G.\ archive; Pietras,
art.\ cit., 179 n.~24, cites the Italian at the same site and quotes it.
The two renderings agree in substance. The recourse's own date is given by
the Commission's answering decree as 13 July 2000.}

\begin{itemize}
\item The Commission received in the rescript \latin{Quia peculiare
munus} of 18 October 1988, at n.~3, the faculty to erect the Fraternity
``according to the particular marks of the Motu Proprio \work{Ecclesia
Dei} n.~6~a, and not merely according to n.~1 of the said rescript, by
which the use of the Roman Missal of 1962 is granted to priests.''
\item Therefore ``the society did not receive an indult after its
creation; it was constituted by reason of that very indult, in order to
realize it communally.'' All members entered freely a society founded for
and by reason of the ancient rite, and thereby accepted, at least
implicitly, to limit the use of their right to celebrate according to the
Missal of common use.
\item The dissensions arose not from personal quarrels but from a
difference of option about that proper characteristic; introducing such a
diversity among the members on a point so essential to their common
spirituality ``will fatally entail the feared disintegration, whoever the
superior chosen.''
\item On term limits: constitution n.~19 does not leave the question
open---it provides that the number of terms is not limited, deliberately so
drafted by the founders in August 1988 and accepted by the Commission in
October. A change would be accepted respectfully but cannot be
retroactive.
\item On the appointment: the 1991 precedent is distinguishable, because
Cardinal Innocenti's decision was taken before the chapter and did not
remove the chapter's right to vote.
\item On the seminaries: no enquiry had been held at the American
seminary, and the canonical visitation of Wigratzbad had underlined the
seriousness and devotion of the teaching body.
\item On the fifth point, the chapter accepts the Roman reading if
``value'' means validity or legitimacy---``it is evident that we adhere to
it fully''---but asks the Commission to state the sense intended, since to
hold the older liturgy of greater value in the ordinary French sense of
the word is to use the right that canon 212 §\,3 recognizes in every one
of the faithful.
\end{itemize}

The recourse ends by proposing, in the alternative, ``a separation of the
society into two institutes, close in their mission but each with its own
proper characteristic, as has happened many times in the history of
institutes of consecrated life.''

\subsection{The decree of 12 October 2000}

The Commission answered by decree, protocol 823/2000, dated 12 October
2000 and signed by Cardinal Castrillón Hoyos. It is short, it is drafted
in the classical form of a rescript with \latin{attento quod} clauses,
and it decides the recourse point by point.\endnote{Pontificia Commissio
\work{Ecclesia Dei}, ``Décret'' 823/2000, 12 October 2000, French text in
the A.M.D.G.\ archive (amdg.asso.fr), read 2026-07-25, expressly answering
the recourse ``déposé au siège de la Commission Pontificale Ecclesia Dei
le 13 juillet 2000'' against five decisions of letter N.~748/2000 of 29
June 2000. Bounded negative: no official publication of this decree was
located, and it is not cited in the canonical article used elsewhere in
this section. It is therefore the least corroborated of the three Roman
acts of 1999--2000 used here, and nothing in this account's conclusions
depends on it alone.}

It \textbf{rejects} points 1 to 3, giving reasons:

\begin{itemize}
\item on term limits: the Commission had already asked for the change,
following a recommendation of the Congregation for Religious concerning
canon 624 §§\,1--2, and the provisional text had been clearly marked as
requiring change in that sense before it could be approved by the Holy
See;
\item on the removal of the chapter's electoral right ``for this
occasion'': it was motivated by a crisis recognized by all and by the
contested composition of the chapter, which failed to ``represent the
whole institute'' or to be ``a true sign of unity in charity'' (canon 631
§\,1), and because the internal government of institutes falls to the
competence of the Holy See (canon 593);
\item on the person appointed: ``the authority to which it belongs to
erect, modify and suppress offices is competent to provide for them''
(canon 148), and the objections raised against the choice of a determinate
person ``are not canonical.''
\end{itemize}

It \textbf{accepts} points 4 and 5: the recourse against the decision on
the seminaries has become moot, the new superior having in fact had the
opportunity to choose the superiors of the houses of formation; and the
fifth request concerns an expression that requires proper interpretation.

\subsection{What was settled, and what was not}

It is easy to read 2000 as a defeat for the Fraternity, and its own
chapter clearly experienced it as one. Read against the documents,
something more precise happened, and the distinction has held for
twenty-six years.

\emph{What Rome refused} was the conversion of the Fraternity's
liturgical patrimony into a personal obligation binding each member and
each seminarian against the universal law. The reason given was not that
the patrimony is unimportant; it is that the general law of the Latin
Church confers on every priest of that Church a faculty which no superior
below the pope can remove, and that a society's particular law may not
achieve by general rule what it may not achieve by precept.

\emph{What Rome did not touch} was the erection, the object stated in the
erection decree, the concession of the 1962 books, or the constitutions'
statement of the institute's own character. On the contrary, the same
letter that refused the special law repeats that the ``holding fast to the
noble traditions in the celebration of holy worship'' is ``the
characteristic mark'' of the Fraternity's particular field in the Church,
and it insists twice that no priest will be forced to use the right it has
defended.

\emph{What Rome asserted} was the reach of canon 593 into internal
government: the power to limit terms of office, to set aside a chapter's
electoral right on a particular occasion, and to appoint a superior
general directly. That is not an extraordinary claim; it is what
pontifical right means when the Apostolic See chooses to exercise it. The
Fraternity's own recourse conceded the principle and disputed only its
application.

\actrecord{Letter of the Pontifical Commission \work{Ecclesia Dei} to the
general chapter, 29 June 2000, Prot.\ N.~748/2000, and decree 823/2000 of
12 October 2000 (German and French texts on partisan archives; the German
letter also cited in the peer-reviewed canonical literature).}{That in
2000 the Holy See refused to make the exclusive use of the 1962 books a
special law of the Fraternity, on the ground that the universal law
confers a right no superior below the pope may remove; limited the
superior general to two consecutive six-year terms; appointed Father
Arnaud Devillers superior general directly, setting aside the chapter's
election on that occasion; required new rectors for both seminaries;
rejected the chapter's recourse on those three heads and accepted it on
two; and did all of this without altering the Fraternity's erection, its
stated object, or its liturgical grant.}{Neither act was published by the
Holy See; both are known through partisan reproductions, the letter
corroborated by its citation in a peer-reviewed canonical article and by
the Fraternity's own contemporaneous documents. The decree of 12 October
2000 is corroborated by nothing but the archive that publishes it. No
claim in this account rests on that decree alone.}
